% 译者说明（Lead 撰写）
\section*{译者说明}
\addcontentsline{toc}{section}{译者说明}

\subsection*{底本与页映射}

本书译自 Lei Fu, \emph{Algebraic Geometry (Second Edition)}（现代数学基础 81，高等教育出版社）的扫描本，共 265 个 PDF 页。原书页码与 PDF 页码的对应关系为
\[
\text{原书页} = \text{PDF 页} - 7
\]
正文中的「原书 p.$N$」一律指原书页码；译注末尾的 \texttt{[依据：原书 p.$N$ 第 $M$ 段]} 即依此定位。

\subsection*{翻译与核对方式}

\begin{enumerate}[label=(\arabic*),leftmargin=2.4em]
  \item 底本为纯图像扫描，无文本层。译文的正文与公式全部对照 200 dpi 页面图像逐句、逐符号核对；OCR 文本仅在个别词形不确定时作参照，公式一律不取自 OCR。
  \item 条目（定义、命题、定理、引理、推论、例、注、构造）沿用原书的单一计数器，编号为「章.节.序号」，与原书一致；习题按节单独编号 1., 2., \ldots。交叉引用使用自动编号，读者可直接与原书对照。
  \item 公式编号按节编排（如 (1.2.1)），与原书一致。
  \item 交换图按原图用 TikZ 重画，箭头方向、标签、对角箭头均对照原图核对。
\end{enumerate}

\subsection*{术语}

术语以「有通行中译则用通行中译，无通行译法或多译并存则保留英文」为原则，全书统一。主要的译法约定如下：sheaf 层、presheaf 前层、section 截面、stalk 茎、germ 芽、scheme 概形、coherent 凝聚、quasi-coherent 拟凝聚、proper morphism 固有态射（不作「真态射」）、separated 分离、quasi-compact 拟紧、flasque sheaf 松弛层、formal completion 形式完全化、spectral sequence 谱序列、derived functor 导出函子、blowing-up 爆破、divisor 除子。人名与专名（Spec、Proj、Ext、Tor、Čech、Cartier、Picard、Zariski、Grothendieck、Serre、Stein、Segre、Koszul、Jacobson、Mittag-Leffler、Hartshorne 等）一律保留拉丁写法，不音译。

\subsection*{译注}

译注以楷体小字印出，标以〔译注〕，只用于三类情况：易错预警、原文含糊或跳步处的点拨、必要的背景直觉。每条译注末尾都注明依据，形式为
\texttt{[来源：\ldots]}（外部文献，附可核查的定位）或 \texttt{[依据：原书 p.$N$ \ldots]}（依据底本自身）。凡无外部依据、仅由上下文推断者，明确标为 \texttt{[译者按（上下文推断，无外部依据）：\ldots]}。本书不引用百科类、问答类网站与个人博客。

\subsection*{关于本译本的说明}

本译本由「数译」工程（编号 SY07-08）以人机协作方式完成：分块翻译、逐块编译自检、全书合稿后统一审计（未定义引用、重复标签、溢出、术语一致性、译注来源完整性）。译文力求忠实于底本，不增删内容；原书笔误一仍其旧，必要时以译注说明。疏漏之处，欢迎读者指正。
